\chapter{Математические формализмы}
\section{Понятия задачи и алгоритма}

В различных математических дисциплинах понятие \textit{задачи} определяется по-разному.
Обычно говорится, что задача - это вопрос, на который можно дать ответ. Определим
данное понятие со стороны теории сложности.

У задачи есть параметры - числовые (множество самих параметров) и "структурные". // TODO: пояснить понятие структурных параметров
В теории сложности задача разбивается на два термина - массовая и индивидуальная задачи.

\begin{definition}
  Массовая задача - это задача с описанным множеством параметров.
\end{definition}

\begin{definition}
  Индивидуальная задача - это задача с фиксированным множеством параметров.
\end{definition}
\sidenote{Пример: задача сложения чисел - массовая задача, а задача сложения чисел 5 и 7 - индивидуальная.}

В теории сложности задачи также различаются по форме, и часто её смена может усложнить задачу.
Обычно задачи разделяют на:
\begin{itemize}
  \item Распознавательные - ответ "да" или "нет"
  \item оптимизационные - ответом является сам объект
  \item вычислительные - ответом является указание параметров объекта
  \item перечислительные - ответом на задачу является число объектов, выполнивших условие
\end{itemize}
\sidenote{Отличие оптимизационной задачи от вычислительной }

На примере одной из самых лёгких задач рассмотрим всевозможные формы.

\begin{abstract}
  Поиск максимального числа
  Заданы N целых чисел: $a_1$, ..., $a_N$. Найдите максимальное из них.
\end{abstract}
Сейчас задача является оптимизационной - и если, например, число $a_7$ является максимальным из всех,
то в качестве ответа нужно выдать число 7, так как оно обозначает данное число, является его индексом.
Переведем вопрос задачи в вычислительную форму - найдите величину максимального числа. Тогда ответом будет само число $a_7$.
С другой стороны, можно переформулировать задачу в форму распознавания - есть ли в данной последовательности
число, большее чем число $B$, где $B$ - константа. Ответом будет являтся да/нет. Перечислительная форма -
найдите количество чисел, больших, чем $B$, где $B$ - константа.

Как видно, изменяя вопрос задачи, можно привести её к любой форме. Рассмотрим наиболее важные задачи,
которые будут встречаться на протяжении всего курса, в их классической форме.

\begin{abstract}
  Разложение на множители
  Дано натуральное число $N$. Требуется представить его как произведение простых множителей.
  Форма: вычислительная.
\end{abstract}

\begin{abstract}
  Простота числа
  Дано натуральное число $N$. Определите, является ли данное число простым.
  Форма: распознавание.
\end{abstract}

\begin{abstract}
  Гамильтонов цикл
  Задан простой неориентированный граф $G = <V, E>$. Определите, есть ли в нём гамильтонов цикл.
  \sidenote{Гамильтонов цикл - это цикл в графе, который проходит через каждую вершину графа ровно 1 раз.}
  Форма: распознавание.
\end{abstract}

\begin{abstract}
  Коммивояжёр
  Задан ориентированный граф $G = <V, E>$, рёбрам которого приписаны веса $c_ij$. Найдите гамильтонов
  цикл экстремальной длины. \sidenote{Экстремальной - или максимальной, или минимальной (чаще всего).}
  Форма: оптимизационная
\end{abstract}

\begin{abstract}
  Задача о выполнимости КНФ\sidenote{КНФ - конъюнктивная нормальная форма - одно из возможных представлений
  булевой функции.}
  Задана КНФ $k(x_1, x_2, ..., x_n)$. Определите, существует ли такой набор $\alpha_1, \alpha_2, ..., \alpha_n$,
  что выполняется $k(\alpha_1, \alpha_2, ..., \alpha_n) = 1$.
  Форма: распознавание.
\end{abstract}

\begin{abstract}
  Задача о клике\sidenote{Клика - подмножество вершин графа, в котором все вершины соединены ребром между собой.}
  Задан граф $G = <V, E>$, необходимо найти максимальную клику.
  Форма: оптимизационная.
\end{abstract}

\begin{abstract}
  Задача о кратчайшем остове\sidenote{Остов - подграф, являющийся деревом.}
  Задан граф $G = <V, E>$, где \(|E| = m\), \(W_1, ..., W_m\) - длины рёбер. Необходимо найти остов наименьшего веса.
  Форма: оптимизационная.
\end{abstract}

\begin{abstract}
  Задача о кратчайшем пути в графе
  Задан граф $G = <V, E>$, где \(|E| = m\), \(W_1, ..., W_m\) - длины рёбер, \(a, b \in V\). Необходимо
  найти минимальный путь между вершинами \(a\) и \(b\).
  Форма: оптимизационная.
\end{abstract}

\begin{abstract}
  Задача об изоморфизме графов. \sidenote{Два графа считаются изоморфными, 
  если можно сопоставить их вершины, не изменяя множество рёбер.}
  Задан графы $G = <V, E>$, $H = <P, Q>$. Изоморфны ли эти графы? 
  Форма: распознавание.
\end{abstract}

Понятие задачи с точки зрения теории сложности определить получилось. Но как для художника - краски,
так и для нас алгоритм - нечто интуитивно понятное, то, что сложно подстроить под какое-то строгое определение.
Если и пытаться это делать, то мы получим одну из вариаций следующего определения - детерминированная
последовательность действий, решающая конкретную задачу. Более строго к этому вопросу подошел Алонзо Чёрч,
представивший следующий тезис (в литературе можно встретить название тезис Чёрча-Тьюринга).

\begin{definition}
  Тезис Чёрча
  Любой алгоритм может быть представлен \textit{конструкцией},
  где конструкция - нормальные алгорифмы Маркова, машина Тьюринга, РАМ, РАСП, язык C++ и так далее.
  \sidenote{Тезис не является истиной в последней инстанции - это предложение об определении
  алгоритма, которым мы успешно пользуемся и его не получается опровергнуть.}
\end{definition}

\section{Нормальные алгорифмы Маркова}

\begin{proposition}
  Любая информация может быть представлена словом в конечном алфавите.
\end{proposition}

Написать про конкатенацию и про пустой символ

В общем смысле любой алгоритм переводит слово \(\alpha\) в алфавите \(A\)
по правилу \(\mathcal{F}\) в слово \(\mathcal{F}(\alpha)\). Более строгое определение
алгорифма Маркова таково:

\begin{definition}
  Алгорифм Маркова - это четвёрка \((A, B, \Lambda, \pi)\), где
  A - полный алфавит,
  B - рабочий алфавит (вспомогательные символы),
  \Lambda - пустой символ
  \pi - конечное множество подстановок
\end{definition}

Отдельно рассмотрим множество \pi и как выглядят подстановки. Простыми словами, подстановка -
это синтаксическая замена в слове одного подслова на другое. Подстановки будем изображать
следующим образом: если \(S_k\) - заменяемое подслово, а \(P_k\) - новое подслово,
то обычная подстановка выглядит как \(S_k -> P_k\). Подстановка, принудительно завершающая
выполнение алгорифма, выглядит так: \(S_k -> (\cdot) P_k\).

При решении задач множество \pi записывается сверху вниз, причем чем выше подстановка, тем
выше её приоритет относительно других в процессе выполнения алгорифма. Типичный процесс
исполнения алгорифма - для входного слова ищется первая подходящая подстановка сверху вниз,
если нашлась, то подстановка применяется и для нового слова снова ищется подстановка сверху вниз.
Если подходящей подстановки больше нет или нашлась подстановка, принудительно завершающая выполнение,
то выполнение алгорифма останавливается.

Приведем пример двух алгорифмов при \(A = {a, b}\):
\begin{equation}
  \begin{orcases}
  \(a -> \Lambda\)
  \(b -> \Lambda\)
  \(\Lambda -> \Lambda\)
  \end{orcases}
\end{equation}
Как видно, этот алгорифм никогда не закончится. А этот уже успешно выполнится:
\begin{equation}
  \begin{orcases}
  \(a -> \Lambda\)
  \(b -> \Lambda\)
  \end{orcases}
\end{equation}

\section{Машина Тьюринга}

\begin{definition}
  Машина Тьюринга - это пятёрка: \((A, S, F, q_{0}, \pi)\), где
  A - алфавит,
  S - множество состояний,
  F - множество конечных состояний, \(F \subseteq S\),
  \(q_{0}\) - начальное состояние,
  \(\pi\) - множество команд
\end{definition}

Одноленточная МТ представляется в виде бесконечной ленты с клетками:
(фото)
Вид команд в Машине Тьюринга:


\begin{definition}
  Машина Тьюринга является алгоритмом, если она останавливается через конечное число тактов.
\end{definition}

(Табличка с незатронутыми ячейками)

Многоленточная МТ состоит из \(m\) лент, за каждый такт пишет и двигается на каждой ленте независимо

Недетерминированная МТ
- язык и задача в форме распознавания
- угадывающая головка
- пишет слово слева, возвращается назад и работает как обычная МТ

Оракульная МТ
- две ленты, обычная и "оракульная"
- оракульная функция, угадывающая головка

\section{Равнодоступная адресная машина}

РАМ (+ задача на *=)

Равномерный вес - в каждой ячейке хранится буква/число
Логарифмический вес - для числа \(a\) \(\lg(a)\) ячеек

РАСП

\section{Кодировки входных данных}

Логарифмический и равномерный веса

\section{Теоремы сравнения}

Почему в тезисе Чёрча именно "конструкция"

НАМ <=> МТ

РАМ <=> РАСП

РАМ к МТ
